
%%%%%%%%%%%%%%%%%%%%%%%%%%%%%%%%%%%%%%%%%%%%%%%%%%%%%%%%%%%%%
\section{模型假设}

为简化问题，本文做出以下假设：

\begin{itemize}[itemindent=2em]
\item \textbf{假设1（$H_1$）：}各电场串联独立工作，总效率为单级效率连乘。ESP多电场标准假设，前级除尘影响后级入口浓度。
\item \textbf{假设2（$H_2$）：}单级除尘效率服从Deutsch方程形式。经典ESP理论方程，物理意义明确，驱进速度$\omega\propto U_i^2$由电晕场与沉积场同电压简化得到。同型号ESP四电场几何相似，设$kA_i=kA_0\cdot g_i$，其中$g=[1.0,1.0,0.9,0.9]$（后电场粉尘更细/比电阻略高，$kA$低约$10\%$），将8参数压缩为5参数缓解欠定。该先验的合理性由AIC信息准则验证：共享$kA_0$模型（5参数）$\Delta\text{AIC}=-9836$强烈优于独立$kA$模型（8参数）。$g_i$的具体取值对结果影响较小——若改为$g=[1,1,0.85,0.85]$或$[1,1,0.95,0.95]$，$kA_0$拟合值相应调整但优化结果变化$<2\%$。
\item \textbf{假设3（$H_3$）：}电耗与电压平方成正比$P_i=k_iU_i^2$，扩展模型加入常数项$c$代表固定空载损耗，约束$k_i\geq 0$保证物理合理性。
\item \textbf{假设4（$H_4$）：}振打周期偏离参考值$T_{ref,i}$双向都使效率下降——周期过长极板积灰抑制沉降，周期过频二次扬尘增加，用双向偏离指数衰减建模，$r=0.5$为过频振打副作用占比。
\item \textbf{假设5（$H_5$）：}工况参数在单个时间窗口内稳定。分钟级数据，K-Means聚类划分典型工况。
\item \textbf{假设6（$H_6$）：}入口条件与操作参数之间无强耦合交互。简化建模，特征重要性验证。
\item \textbf{假设7（$H_7$）：}振打瞬时峰值仅与振打周期和当前浓度相关。半经验关系，分钟级数据无法直接观测秒级峰值。
\item \textbf{假设8（$H_8$）：}电除尘器工作在电晕放电伏安特性单调区，二次电压与二次电流一一对应，调节二次电压等价于调节二次电流，不单独建模电流（电耗$U^2$近似$U\cdot I$，隐含$I\propto U$）。
\end{itemize}

\textbf{物理先验汇总与可靠性保障。}上述假设中包含若干无法从限幅数据中辨识的参数，其取值依赖物理先验或工程经验。本文对这些先验的可靠性采取三重保障策略：（1）\textbf{信息准则验证}——共享$kA_0$（5参数）vs独立$kA$（8参数）由AIC判定（$\Delta\text{AIC}=-9836$），避免过参数化；（2）\textbf{敏感性分析}——$r\in[0.3,1.0]$时$P^*$几乎不变，$g_i$变化$\pm0.05$时优化结果变化$<2\%$，$kA_0$在贝叶斯95\% CI内电耗变化$<1\%$；（3）\textbf{不可辨识性披露}——$r$、$T_{ref}$、$\alpha_{3,4}$的不可辨识性在6.3节和5.2.2节明确讨论，并建议通过后续实验标定。对于峰值估算中的经验系数0.1，因分钟级数据无法验证，论文将其定位为量级估算而非精确结论（详见5.1.2节步骤5）。
